\documentclass[11pt]{article}
\usepackage[T1]{fontenc}
\usepackage[margin=1in]{geometry}
\usepackage{amsmath,amssymb}
\usepackage{enumitem}
\usepackage{hyperref}
\hypersetup{colorlinks=true,linkcolor=black,urlcolor=blue}

\begin{document}

\begin{center}
{\LARGE Breathing Dynamics Formula Sheet}\\[0.5em]
{\large PolarH10 breath interval and amplitude entropy features}
\end{center}

\section*{Data path}

\begin{itemize}[leftmargin=1.5em]
\item input waveform: calibrated \texttt{VolumeBase01} from the Breathing tracker
\item event extraction: alternating peaks and troughs from that calibrated waveform
\item derived series:
  \begin{itemize}[leftmargin=1.5em]
  \item interval series: spacing between same-polarity extrema
  \item amplitude series: absolute peak-trough excursion
  \end{itemize}
\end{itemize}

This is the main adaptation relative to the cited paper: the paper uses a
respiration-belt signal, while this app uses the repository's calibrated ACC
breathing waveform.

\section*{Derived breath series}

For accepted extrema times $t_k$ and values $x_k$:

\[
\text{interval}_j = t_j - t_{j-1}
\]

where the two extrema belong to the same polarity family, for example peak to
peak or trough to trough.

\[
\text{amplitude}_j = |x_j - x_{j-1}|
\]

where the two extrema are alternating peak and trough partners in a completed
breath excursion.

\section*{Basic statistics}

For either derived series $s_1 \ldots s_n$:

\[
\mu = \frac{1}{n}\sum_{i=1}^{n}s_i
\]

\[
SD = \sqrt{\frac{1}{n-1}\sum_{i=1}^{n}(s_i - \mu)^2}
\]

\[
CV = \frac{SD}{|\mu|}
\]

\section*{ACW50}

The autocorrelation window is the first lag where the normalized
autocorrelation falls below $0.5$:

\[
ACW50 = \min \left\{ \ell \ge 1 : \rho(\ell) < 0.5 \right\}
\]

\section*{PSD slope}

The implementation now follows the cited NeuroKit2 path more closely:

\begin{itemize}[leftmargin=1.5em]
\item linearly detrend the series
\item standardize it with population SD
\item compute a one-sided FFT power spectrum
\item keep the lower-frequency quartile of the spectrum
\item fit the slope in log-log space
\end{itemize}

\[
\text{PSD slope} =
\operatorname{slope}\left(\log_{10}(f), \log_{10}(PSD(f))\right)
\]

\section*{Lempel-Ziv complexity}

The series is binarized around its mean:

\[
b_i =
\begin{cases}
1 & s_i \ge \mu \\
0 & s_i < \mu
\end{cases}
\]

The Lempel-Ziv counter then scans the binary string left-to-right and adds a new
phrase whenever it encounters a substring not yet present in the discovered
pattern set. The app publishes a normalized form:

\[
LZC_{norm} = \frac{c(n)}{n / \log_2(n)}
\]

\section*{Sample entropy}

Default settings:

\begin{itemize}[leftmargin=1.5em]
\item $m = 2$
\item delay $= 1$
\item $r = 0.2 \cdot SD$
\end{itemize}

\[
SampEn(m, r, N) = -\ln\left(\frac{A}{B}\right)
\]

where $B$ counts matching embedded vectors of length $m$, $A$ counts matching
embedded vectors of length $m + 1$, and matching uses the Chebyshev-distance
tolerance $r$.

If the conditional probability is undefined, the app now suppresses the
headline instead of inventing a finite fallback value.

\section*{Multiscale entropy}

Default settings:

\begin{itemize}[leftmargin=1.5em]
\item $m = 3$
\item delay $= 1$
\item $r = 0.2 \cdot SD$ of the original series
\item scales $1..5$
\item published summary: area under the entropy-by-scale curve divided by the count of usable scales
\end{itemize}

For each scale $\tau$, the series is coarse-grained by averaging
non-overlapping blocks:

\[
y_j^{(\tau)} = \frac{1}{\tau}\sum_{i=(j-1)\tau + 1}^{j\tau} s_i
\]

Sample entropy is then computed on each coarse-grained series with the same
base tolerance. The reported MSE summary is:

\[
MSE_{AUC} = \frac{\operatorname{trapz}(SampEn_{\tau})}{K}
\]

where $K$ is the number of usable finite scales kept for the summary.

\section*{Alignment status}

This implementation is aligned to the cited method family for interval and
amplitude series extraction from extrema; mean, SD, CV, ACW50, PSD slope, LZC,
SampEn, and MSE; NeuroKit2-style SampEn and MSE defaults; NeuroKit2-style
handling of undefined entropy values; and NeuroKit2-style PSD slope
preprocessing and low-frequency fit region.

The remaining explicit adaptation is the waveform source: the paper uses
respiration-belt input, while this repository uses the calibrated ACC breathing
waveform from \texttt{PolarBreathingTracker}.

\section*{References}

\begin{itemize}[leftmargin=1.5em]
\item Goheen et al., ``It's About Time: Breathing Dynamics Modulate Emotion and Cognition'' (2025).
\item \url{https://doi.org/10.1111/psyp.70149}
\item \url{https://mesmerprism.github.io/PolarH10/reference/breathing-dynamics-workflow.html}
\item \url{https://mesmerprism.github.io/PolarH10/reference/references.html}
\end{itemize}

\end{document}
